\chapter*{Lời Mở Đầu}
\addcontentsline{toc}{chapter}{Lời Mở Đầu}
\markboth{Lời Mở Đầu}{Lời Mở Đầu}

\begin{truyen}{Lời Mở Đầu}{Opening Words}
\chuhoa{L}{ớp học không chỉ có bài giảng và điểm số. Nó còn có những câu hỏi nghe quen đến mức giáo viên chưa cần nghe hết đã biết học sinh đang muốn né gì, sợ gì, hay đang chống chế cho điều gì.}
\textit{(A classroom is not just lessons and grades. It is also full of familiar questions that already reveal what the student is avoiding, fearing, or trying to excuse.)}

Cuốn sách này gom lại 1000 câu hỏi kiểu đó và cho chúng một câu đáp đủ duyên để bật cười, đủ thâm để nhớ lâu, và đủ nhẹ để không biến lớp học thành nơi mỉa mai độc địa.
\textit{(This book gathers 1000 such questions and gives each one a reply that is funny enough to spark laughter, sharp enough to stick, and light enough to avoid cruelty.)}

Đọc cuốn này, người ta sẽ thấy một điều rất quen: học sinh hỏi câu nào cũng tưởng mới, nhưng thật ra giáo viên đã nghe các phiên bản của nó suốt nhiều năm. Cũng vì vậy mà một câu đáp hay đôi khi quý hơn một bài giảng dài.
\textit{(Readers will notice something familiar: students often think their excuse is new, but teachers have heard versions of it for years. That is why one good reply can sometimes matter more than a long lecture.)}
\end{truyen}

\begin{baihoc}
Một câu trả lời duyên có thể làm học sinh cười. Một câu trả lời đúng có thể làm học sinh nhớ rất lâu. Khi hai thứ đi chung, lớp học sẽ có sức sống hơn nhiều.
\textit{(A witty answer may make students laugh. A true answer may stay with them for years. When both come together, the classroom becomes far more alive.)}
\end{baihoc}

\ngancach